\documentclass[11pt]{article}

\usepackage[margin=1in]{geometry}
\usepackage{booktabs}
\usepackage{hyperref}

\title{Learning Signal Density: Measuring External-Sample and Internal-Compute Efficiency in Learning Loops}
\author{Francesco Giannicola \\ Limes Labs}
\date{2026}

\begin{document}
\maketitle

\begin{abstract}
This paper studies whether apparent sample efficiency is better attributed to a
whole learning loop than to a neural network in isolation. We define external
sample efficiency, compute efficiency, and learning-signal density, then use
controlled synthetic domains to compare raw training examples with selected,
transformed, replayed, and feedback-enriched training loops. The current
repository contains a pilot audit and a preregistered path toward tiny neural
replications. The pilot is not a claim about frontier language models.
\end{abstract}

\section{Motivation}

Comparisons between human learning and language-model pretraining often count
raw tokens or approximate brain compute. This can hide the central mechanism:
biological and engineered systems do not merely receive data. They select,
transform, replay, reward, and consolidate it.

\section{Definitions}

Let a learning system be
\[
S = (\theta, M, \pi_{\mathrm{collect}}, \pi_{\mathrm{select}},
\pi_{\mathrm{transform}}, \pi_{\mathrm{train}}, R, E).
\]

We measure:
\[
\mathrm{external\ sample\ efficiency} =
\frac{\Delta Q}{N_{\mathrm{external\ events}}},
\]
\[
\mathrm{compute\ efficiency} =
\frac{\Delta Q}{C_{\mathrm{internal}}},
\]
and
\[
\mathrm{learning\ signal\ density} =
\frac{\Delta Q}{N_{\mathrm{external\ events}} C_{\mathrm{internal}}}.
\]

\section{Pilot Experiment}

The first experiment uses a synthetic causal-text domain with hidden rules,
fixed train/validation/heldout splits, and several train-only processing
pipelines. It records heldout accuracy, majority-baseline improvement,
external-event count, internal-token count, charged transform costs, signed
efficiency metrics, and sample-budget sensitivity. The first non-oracle
transformation condition induces empirical rules from the train split before
generating counterfactual labels. A validation-gated variant selects induction
thresholds on validation data and charges the tuning overhead.
A direct validation-gated variant selects thresholds by induced-rule
precision/coverage without retraining the learner for every candidate.
An MDL-style variant selects compact empirical rules with a description-length
penalty before generating counterfactuals.

\section{Limitations}

The current pilot uses an online linear learner, not a neural language model.
Counterfactual labels are oracle-generated inside the synthetic world. These
limitations are intentional for the first audit and must be removed or bracketed
before a stronger paper claim.

\bibliographystyle{plain}
\bibliography{references}

\end{document}
